\subsection{The four phases of the proof}
\label{sec:four-proof-phases}

A \emph{role} is a vertex of the cube being embedded, and a \emph{label}
is a vertex of the host graph. We map the even and odd roles to disjoint
host sides. The construction must preserve every cube edge in one color
and assign distinct labels to distinct roles. After an injective assignment
$f$ of the odd roles, the possible labels for an even role $v$ form the
common-neighbor set
\[
 L_v=\bigcap_{u\sim v}N_G(f(u)),
\]
where $G$ is the chosen color. Most embedding arguments below construct
probability laws $p_v$ supported on $L_v$ with
$\sum_v p_v(x)\le1$ for every host label $x$. Summing over any set of roles
then gives Hall's condition and hence distinct representatives. In the
last part of the proof, we instead construct random two-element subsets
of $L_v$ and bound directly the probability that any collection of these
sets violates Hall's condition.

To describe the reductions, a \emph{patch} will mean a pair of host sets,
one on each side, with specified sampling laws $\mu,\nu$; the required
density properties vary between the constructions. Write
\[
 d_G(\mu,\nu)=\Pr\{xy\in G\},\qquad x\sim\mu,\quad y\sim\nu
 \quad\text{independently}.
\]
The \emph{width} of $\mu$ is $\log(N\max_x\mu(x))$: a width bound $w$
allows no atom larger than $e^w/N$. In particular, the uniform law on
$M$ labels has width $\log(N/M)$. A discrepancy bound controls
$|d_G(\mu,\nu)-1/2|$ for every pair of laws within specified width budgets.
An orientation specifies which host side receives the first law; reversing
it interchanges the sides and their width budgets.

A patch property is \emph{available} if, for some fixed $\kappa>0$ and
all sufficiently large $n$, a patch satisfying it exists after every
removal of at most $\kappa N$ labels from each side. The patch may
change with the removed sets.

\begin{enumerate}[leftmargin=*]
\item \textbf{Reduce to a sequence and construct the probability tools.}
If the theorem fails, Section~\ref{sec:reduction-and-notation} gives
counterexamples with two disjoint host sides of size $N$, where
$n\to\infty$, $N/2^n\to\infty$, and $N\le n2^n$. The construction must
work without any prescribed rate for the middle limit.
Section~\ref{sec:general-conventions-and-probabilistic-tools} provides
mixtures of patch laws with small average label masses, locally determined
choices of random data, and ways to sample labels injectively while
controlling their joint probabilities.

\item \textbf{Use density gains or deduce discrepancy.}
By Lemma~\ref{lem:stabilization}, we may pass to a subsequence and
remove $o(N)$ labels from each side so that each tested patch property is
either absent or available. At prescribed power widths,
Section~\ref{sec:bias-versus-purity-at-power-widths} then shows that
available bias of at least a suitable negative power of $n$ forces an
available nearly monochromatic patch at doubled widths, unless a cube can
already be embedded. The broad-side and
small-grid constructions in
Sections~\ref{sec:a-constant-broad-high-neighborhood-case}--\ref{sec:a-small-power-key-grid}
then embed a cube under their respective concentration and high-degree
hypotheses. These implications, applied in both colors of a counterexample,
exclude bias throughout an initial range of widths: every permitted law
pair has density close to one half.
Sections~\ref{sec:asymmetric-pure-patches-under-2}--\ref{sec:jump-to-large-bias-only-at-linear-budget}
extend this range, allowing much greater concentration on one side than
on the other. They also exclude the full-dimensional cluster witnesses
of Section~\ref{sec:full-dimensional-cluster-patches}, in which many
labels have unusually large pairwise common neighborhoods.
Section~\ref{sec:the-two-deep-orientations-some-estimates} records the
resulting discrepancy bounds with a small linear width budget on either
side and a power width on the other. The cluster exclusion also bounds
the scales of the smaller cluster patches used next.

\item \textbf{Allocate the cube to patches and use the larger gains.}
Section~\ref{sec:patch-tiling-under-9} extracts disjoint patches and
retains a family with one color, one orientation, and one of the finitely
many constructions described there. It allocates the cube to these
patches by subcubes obtained by fixing initial coordinates, with sizes
proportional to the patch sizes. This is the geometry used in earlier
cube-versus-clique tilings~\cite{conlon-fox-lee-sudakov2016,fiz-pontiveros-griffiths-morris-saxton-skokan2016}.
The fixed coordinates identify the cube edges crossing between patches.

Smaller host sets and crossing edges reduce the available common
neighborhoods, so the construction needs a compensating gain. For
comparison, suppose that each of $n$ independent restrictions retains
each label with probability $1/2+\delta$. A set of size $M$ would then
leave an expected $M(1/2+\delta)^n$ labels. Relative to $N2^{-n}$, the factor is
$(M/N)(1+2\delta)^n$: the accumulated density surplus can compensate for
the smaller support. The proof obtains the needed gain either from edge
density or from common-neighbor density inside cluster patches, and
bounds the additional losses from crossing edges and dependent choices.
Section~\ref{sec:internal-solver-and-profiles-for-these-tilings}
constructs the laws used inside cluster patches. When the gain is large
enough, Section~\ref{sec:high-modes-of-the-tiling} assigns the odd roles
injectively and bounds the column sums of the even laws supported on
$L_v$. Hall's theorem then completes the embedding.

\item \textbf{Complete the cases with smaller gains.}
Sections~\ref{sec:low-modes-geometry-and-cell-states}
and~\ref{sec:initial-list-gates-and-resampling-estimates} first assign
the odd roles and preserve lists for the even roles. A fixed linear map
on cube words partitions the odd roles into syndrome classes; a chosen
$O(\log n)$ of these classes receive provisional labels. Their labels
are replaced later, one class at a time, from reserved host sets. Each
even role has at most one neighbor in each late class.
Section~\ref{sec:late-dynamics-and-its-transfer-problem} controls how
these replacements shorten its list and leaves two distinct candidates
adjacent to all its assigned odd neighbors. A joint bound on the
candidate pairs rules out Hall-deficient collections, giving distinct
even labels.
\end{enumerate}

Appendix~\ref{app:regime-table} collects the quantitative branches,
including their width, density, and loss budgets. It is a reference for
the parameter choices after the corresponding constructions have been
introduced.

\paragraph{The probability comparisons.}
A recurring construction samples a hidden $k$-tuple on the first host
side and requires neighboring odd labels to be adjacent to every entry.
Given these odd labels, the tuple's conditional distribution is supported
on the $k$th Cartesian power of their common neighborhood. A bound on its largest atom
therefore gives a lower bound on the size of that neighborhood. The tuple
can also influence which observations are selected. Its likelihood must
include that selection event, with the selection rule recomputed when the
tuple is varied. Lemma~\ref{lem:gated-posterior} formulates this comparison;
Section~\ref{sec:bias-versus-purity-at-power-widths} gives its first
application.

Injectivity introduces another dependence. The calibrated sampler in
Lemma~\ref{lem:near-product-injection} preserves individual marginals
exactly and bounds probabilities of several specified outputs by their
product probabilities times a controlled factor. Exact marginals are
needed in later cluster constructions whose scales may remain fixed.
The clock sampler in Lemma~\ref{lem:clock-sampling} also avoids specified
local failures and provides a joint upper bound. Summing such bounds
against nonnegative functions permits the common-neighbor and load
estimates to be used after injective sampling. The later sequential
assignment requires further conditional comparisons, developed where the
successive observations are defined in
Section~\ref{sec:late-dynamics-and-its-transfer-problem}.
